\documentclass{article}
\usepackage[T1]{fontenc}
\usepackage{lmodern}
\usepackage{graphicx}
\usepackage{amsmath,amssymb}
\usepackage[a4paper,margin=2cm]{geometry}
\usepackage[absolute,overlay]{textpos}
\setlength{\TPHorizModule}{1mm}
\setlength{\TPVertModule}{1mm}
\usepackage[indonesian]{babel}
\usepackage{hyperref}
\setlength{\emergencystretch}{2em}
% unit: Halaman 13

\begin{document}

% === Halaman 13 (python/native) ===

AES-128

𝐾0

𝐾1

𝐾2

𝐾9

(1) SubBytes

(2) ShiftRow

(3) MixColumn

⊕

(1) SubBytes

(2) ShiftRow

(3) MixColumn

(1) SubBytes

(2) ShiftRow

⊕

⊕

⊕

⊕

𝐾10

input

output

⋯

10 rounds

𝐾

128

128

128

128

128

Sumber: Håkon Jacobsen, Lecture 2 – Block ciphers, PRFs/PRPs, DES, AES

\end{document}
